\subsection{Spectral Metrics}
\label{subsec:spectral_metrics}

Given a language model at generation $t$, we extract last-layer embeddings for $n{=}1{,}000$ held-out samples to form $\mathbf{E} \in \mathbb{R}^{n \times d}$ and compute the Gram matrix $\mathbf{G} = \mathbf{E}^\top \mathbf{E}$.
Let $\sigma_1 \geq \sigma_2 \geq \cdots \geq \sigma_d$ denote the singular values of $\mathbf{E}$.
We track three complementary spectral indicators:

\paragraph{Spectral Log-Volume (SLV).}
$\text{SLV} = \sum_i \log \sigma_i = \frac{1}{2} \log \det \mathbf{G}$.
This quantity measures the total volume spanned by the embedding cloud in log-space.
Because it sums over all singular values, SLV is sensitive to contraction along any spectral direction, not only the dominant ones.

\paragraph{Effective Rank.}
$r_{\text{eff}} = \exp(H(\mathbf{p}))$, where $p_i = \sigma_i / \sum_j \sigma_j$ and $H$ is the Shannon entropy.
Effective rank quantifies how many dimensions carry substantial variance.
A uniform spectrum yields $r_{\text{eff}} \approx d$; concentration onto few directions drives $r_{\text{eff}}$ toward~1.

\paragraph{Leading Singular Value.}
$\sigma_1$ tracks the energy in the most dominant spectral direction.
Rising $\sigma_1$ under falling $r_{\text{eff}}$ signals that variance is concentrating into a low-dimensional subspace, a hallmark of representational collapse.

These three metrics capture complementary collapse signatures: SLV detects global volume loss, $r_{\text{eff}}$ measures dimensionality reduction, and $\sigma_1$ isolates spectral concentration.

\subsection{Adaptive Mixing Controller}
\label{subsec:controller}

We propose a closed-loop controller that adjusts the real-data mixing ratio $r_t$ at each generation based on the rate of spectral change.
The update rule is:
\begin{equation}
\label{eq:controller}
r_{t+1} = \frac{B_{\text{rem}}}{T - t} + k \cdot \max(0,\; \Delta_t - \tau)
\end{equation}
where $B_{\text{rem}}$ is the remaining real-data budget (fraction of total budget not yet consumed), $T$ is the total number of generations, $\Delta_t$ is the magnitude of the per-generation SLV change, $\tau{=}0.002$ is a sensitivity threshold, and $k{=}10$ is the controller gain.
The output is clipped to $[r_{\min}, r_{\max}] = [0.20, 0.80]$ and subject to the global constraint $\sum_t r_t \leq B_{\text{total}}$.

The first term distributes the remaining budget uniformly across future generations, ensuring graceful degradation even without spectral signal.
The second term injects additional real data when SLV contraction exceeds the threshold $\tau$, proportional to the severity of the spectral drop.
This design requires no lookahead and operates with a single hyperparameter choice ($k$) beyond the budget constraint.

\subsection{Experimental Conditions}
\label{subsec:conditions}

We compare four mixing strategies that share the same total real-data budget to isolate the effect of allocation timing:

\begin{enumerate}
\item \textbf{Pure collapse} ($r_t{=}0$): all training data is synthetic after Gen~0. This baseline characterizes unmitigated collapse.
\item \textbf{Fixed mix} ($r_t{=}0.50$): a static 50/50 real-synthetic ratio at every generation.
\item \textbf{Linear increasing}: $r_t$ grows linearly from 0.20 to 0.80 across generations, an open-loop heuristic that front-loads synthetic data.
\item \textbf{Adaptive} ($k{=}10$): the SLV-guided controller from Eq.~\ref{eq:controller}.
\end{enumerate}

All conditions use Pythia-410M-deduped fine-tuned on a 50K-sample OpenWebText subset for 10 recursive generations.
We evaluate with perplexity, HellaSwag accuracy (0-shot), distinct-4, and the three spectral metrics defined above.